\documentclass[11pt]{article}
\usepackage{fontspec}
\setmainfont{Times New Roman}
\setsansfont{Arial}
\setmonofont{Consolas}
\usepackage{amsmath,amssymb,mathtools}
\usepackage{geometry}
\usepackage{microtype}
\newcommand{\mS}{\mathfrak{S}}
\newcommand{\mA}{\mathfrak{A}}
\geometry{a4paper,margin=2.35cm}
\setlength{\parindent}{1.25em}
\setlength{\parskip}{0pt}
\makeatletter
\renewcommand\footnotesize{\@setfontsize\footnotesize{7.6}{8.6}\selectfont}
\makeatother
\emergencystretch=100em
\allowdisplaybreaks
\sloppy
\hbadness=10000
\vbadness=10000
\hfuzz=2pt
\vfuzz=10pt
\raggedbottom

\begin{document}

% Unit: Paper32 section 2 quaternion idempotent v001. Direct Interslavic Latin translation from the R122 German Paper32/Paper33 source slice, lines 395--420 / segment P32-S0014.
\noindent Za tu neograničenost nyně kako primjer vozměmo kvaternionno tělo, vzeto kako hyperkompleksny sistem v odnosu k polju $P$ racionalnyh čisel; s bazisnymi elementami $i,j,k$ kromě jedinice $1$ i s poznatymi relacijami:
\[
 i^2=j^2=k^2=-1;\qquad ij=k;\quad jk=i;\quad ki=j;\quad ji=-k;\quad kj=-i;\quad ik=-j.
\]
Pokáže se, že razpadne polja sut vsě i samo te algebraične čislove polja, pri pridruženju ktoryh obstaje idempotentny element $r$, to jest takovy, že $r^2=r$, pri čem $r$ jest različny i od nuljevogo, i od jediničnogo elementa. Ta čislove polja možno takože harakterizovati kako te, v ktoryh $(-1)$ jest predstavimo kako suma trěh kvadratov, a slědovatelno i kako suma dvuh kvadratov.\footnote{Že iz predstavimosti $(-1)$ kako sumy trěh kvadratov vsegda slěduje predstavimost kako sumy dvuh kvadratov, pokazuje identičnost
\[
(c^2+d^2)(a^2+b^2+c^2+d^2)=(ac+bd)^2+(ad-bc)^2+(c^2+d^2)^2,
\]
kotra, napr., slěduje iz teorema o produktu norm kvaternionov. Naměsto zahtěvati obstajanje idempotentnogo elementa, za razpad bylo by dostatočno zahtěvati element s nulevoju normoju.} Vpravdu, ako položimo
\[
        r=\alpha+\frac12\beta i+\frac12\gamma j+\frac12\delta k,
\]
to
\[
        r^2=\left(\alpha^2-\frac14\beta^2-\frac14\gamma^2-\frac14\delta^2\right)
          +\alpha\beta i+\alpha\gamma j+\alpha\delta k,
\]
i zato uslovje $r^2=r$ jest ravnoznačno sistemu
\[
4\alpha^2-4\alpha-\beta^2-\gamma^2-\delta^2=0;\quad
2\alpha\beta-\beta=0;\quad 2\alpha\gamma-\gamma=0;\quad 2\alpha\delta-\delta=0,
\]
a tako, poneže $\beta,\gamma,\delta$ ne smut izčeznuti jednovremenno, ravnoznačno uslovju
\[
        (-1)=\beta^2+\gamma^2+\delta^2.
\]

\end{document}
